\section*{Chapter \ref{ch:b2:pt-paths} --- Thermobarometry and P--T--t Paths}

\begin{solution}{exo:b2:pt-paths:1}
(a) A chlorite zone, B biotite zone, C garnet zone, D staurolite zone, E kyanite zone,
F sillimanite zone. (b) Five: the biotite, garnet, staurolite, kyanite and sillimanite
isograds, each between two successive samples, grade rising from A to F. (c) An isograd
only says that a reaction ran; a field gradient needs pressures and temperatures, which
come from calibrated reactions or from thermobarometry of the peak assemblage of each
rock.
\end{solution}

\begin{solution}{exo:b2:pt-paths:2}
(a) $(X_{\text{Mg}}/X_{\text{Fe}})_{\text{grt}} = 0.10/0.80 = 0.125$,
$(X_{\text{Mg}}/X_{\text{Fe}})_{\text{bt}} = 0.45/0.50 = 0.900$; $K_D = 0.1389$, $\ln K_D =
-1.9741$. (b) $3R = 24.942$; denominator $20.19 + 24.942\times1.9741 = 69.428$; at
\qty{5}{kbar}, $T = 54185/69.428 = \qty{780.5}{K} = \qty{507}{\celsius}$; at
\qty{7}{kbar}, $T = (52630 + 2177)/69.428 = \qty{789.4}{K} = \qty{516}{\celsius}$.
(c) Two kilobars change it by only \qty{9}{K}: a thermometer.
\end{solution}

\begin{solution}{exo:b2:pt-paths:3}
(a) $z = P/(\bar\rho g) = 6.0\times10^8/(2800\times9.81) = \qty{21.8}{km}$.
(b) $(580 - 10)/21.8 = \qty{26}{\celsius\per\km}$. (c) At \qty{580}{\celsius} the line
lies at $3.67 + 0.8\times2.10 = \qty{5.35}{kbar}$; at \qty{6.0}{kbar} the rock is on the
high-pressure side: kyanite.
\end{solution}

\begin{solution}{exo:b2:pt-paths:4}
(a) Reactants $2\times140.693 + 3\times18.015 = \qty{335.431}{g}$; products $277.112 +
58.320 = \qty{335.432}{g}$: balanced to the rounding. (b) $2\times43.65 =
\qty{87.30}{cm^3}$ become $107.50 + 24.63 = \qty{132.13}{cm^3}$: $+\qty{51.4}{\percent}$.
(c) Olivine $281.386/87.30 = \qty{3.223}{g.cm^{-3}}$; products $335.432/132.13 =
\qty{2.539}{g.cm^{-3}}$. (d) The rock swells by half its volume and becomes lighter; it
also stores water in its minerals.
\end{solution}

\begin{solution}{exo:b2:pt-paths:5}
(a) Means $-1.206$ at \qty{1072.15}{K} and $-1.782$ at \qty{823.15}{K}. With $1/T =
9.3271\times10^{-4}$ and $1.21485\times10^{-3}$, $a = (-1.206 + 1.782)/(9.3271\times10^{-4}
- 1.21485\times10^{-3}) = -2041.5$ K and $b = -1.206 + 2041.5/1072.15 = 0.6982$.
(b) $\Delta_r H + 2070\times0.311 = 2041.5\times24.942 = \qty{50920}{J}$, so $\Delta_r H =
\qty{50276}{J}$; $\Delta_r S = 0.6982\times24.942 = \qty{17.41}{J.K^{-1}}$. (c) At
\qty{5}{kbar}: for $K_D = 0.18$, 588.1 against \qty{587.5}{\celsius}; for $0.10$, 419
against \qty{425}{\celsius}; for $0.30$, 819 against \qty{806}{\celsius}. Within the range
of the experiments the two calibrations agree, though their constants differ by
\qty{2.4}{kJ} and \qty{2.8}{J.K^{-1}}; the errors of $\Delta_r H$ and $\Delta_r S$
compensate each other there and grow outside it. A calibration should not be used far
outside the conditions of its experiments.
\end{solution}

\begin{solution}{exo:b2:pt-paths:6}
(a) Exchange: $(113.12 + 154.30) - (115.32 + 149.65) = \qty{+2.45}{cm^3} =
\qty{0.245}{J.bar^{-1}}$, against 0.311. GASP: $125.28 + 2\times44.15 + 22.69 -
3\times100.79 = \qty{-66.10}{cm^3} = \qty{-6.61}{J.bar^{-1}}$, against $-6.21$, the value
of the calibration at \qty{800}{\celsius} and \qty{10}{kbar}, where thermal expansion and
compression change the volumes. In both cases the GASP change is twenty to thirty times
the exchange one. (b) $\partial T/\partial P = 0.311/62.96 = \qty{4.9e-3}{K/bar}$:
\qty{4.9}{K} per kilobar. (c) The slope of the line is $22.80 + 8.314\times3\ln 0.20/6.21
= 22.80 - 6.46 = \qty{16.3}{bar/K}$; \qty{17}{K} shift the pressure by
\qty{280}{bar}.
\end{solution}

\begin{solution}{exo:b2:pt-paths:7}
(a) $P_0 = 22.80\times600 - 1093 = \qty{12587}{bar}$; $3\ln(0.08/0.30) = -3.9653$; $RT =
8.314\times873.15 = \qty{7259.4}{J}$; $P = 12587 + 7259.4\times(-3.9653)/6.21 = 12587 -
4635 = \qty{7952}{bar}$. (b) $3\ln(0.08/0.35) = -4.4277$, $P = 12587 - 5176 =
\qty{7411}{bar}$: \qty{0.54}{kbar} less, so the plagioclase must be analysed with care.
(c) $7952\times10^5/(2800\times9.81) = \qty{29.0}{km}$.
\end{solution}

\begin{solution}{exo:b2:pt-paths:8}
(a) $z = P/(\bar\rho g)$: \qty{118.9}{km}, \qty{34.0}{km} and \qty{15.3}{km}.
(b) $(118.9 - 34.0)/2.2 = \qty{38.6}{km/Myr} = \qty{3.9}{cm.yr^{-1}}$ and $(34.0 -
15.3)/1.1 = \qty{17.0}{km/Myr} = \qty{1.7}{cm.yr^{-1}}$, a little above the published
3.4 and 1.6: their conversion of pressure to depth gives slightly smaller depth
differences. (c) $\sigma = \sqrt{0.9^2 + 0.9^2} = \qty{1.3}{Myr}$, so $2.2 \pm 1.3$
million years and a rate between $84.9/3.47 = 24$ and $84.9/0.93 = \qty{91}{km/Myr}$,
2.4 to \qty{9.1}{cm.yr^{-1}}: fast in any case. (d) Even the lowest rate is about fifty times
the \qty{0.5}{mm.yr^{-1}} of the model: erosion alone cannot do it.
\end{solution}

\begin{solution}{exo:b2:pt-paths:9}
(a) $\partial T/\partial\ln K_D = 3RT/(\Delta_r S - 3R\ln K_D) = 24.942\times860.6/62.96 =
\qty{341}{K}$; $\partial T/\partial P = 0.311/62.96 = \qty{4.9}{K}$ per kilobar.
(b) $\sigma_T = \sqrt{(341\times0.05)^2 + 4.9^2} = \sqrt{17.0^2 + 4.9^2} =
\qty{17.7}{K}$. (c) The analytical error on $K_D$ dominates; it alone is two thirds of
the \qty{25}{\celsius} of the best calibrations, which also include the errors of the
calibration itself.
\end{solution}

\begin{solution}{exo:b2:pt-paths:10}
(a) Garnet $0.10/0.75 = 0.1333$, biotite $0.42/0.52 = 0.8077$, $K_D = 0.1650$, $\ln K_D =
-1.8018$, denominator $20.19 + 24.942\times1.8018 = 65.130$; $3\ln(0.07/0.25) = -3.8189$.
At \qty{6000}{bar}: $T = 54496/65.130 = \qty{836.7}{K} = \qty{563.6}{\celsius}$, then $P
= 22.80\times563.6 - 1093 + 8.314\times836.7\times(-3.8189)/6.21 = \qty{7479}{bar}$.
At \qty{7479}{bar}: $T = \qty{570.6}{\celsius}$, $P = \qty{7603}{bar}$. At
\qty{7603}{bar}: $T = \qty{571.2}{\celsius}$, $P = \qty{7614}{bar}$; once more, $T =
\qty{571.3}{\celsius}$. The rock equilibrated at about \qty{571}{\celsius} and
\qty{7.6}{kbar}. (b) At \qty{571}{\celsius} the kyanite $=$ sillimanite line lies at $3.67
+ 0.71\times2.10 = \qty{5.16}{kbar}$: kyanite. (c) The barometer was calibrated for
kyanite; with sillimanite the reaction, and its line, differ, since the two polymorphs do
not have the same Gibbs energy.
\end{solution}

\begin{solution}{exo:b2:pt-paths:11}
(a) $\tau = L^2/(\pi^2\kappa)$: for \qty{35}{km}, $1.225\times10^9/(9.870\times10^{-6}) =
\qty{1.24e14}{s} = \qty{3.9}{Myr}$; for \qty{70}{km}, four times more,
\qty{15.7}{Myr}. (b) Seventy million years is four to eighteen times these times: heat
moves through the thickened crust much faster than erosion removes it. (c) The rock
reaches its greatest depth at once, in the stacking; heat arrives over the following
tens of millions of years, while erosion is slowly bringing it up: the temperature keeps
rising for a while during the decompression. (d) Erosion in seven million years would
bring the rock up before conduction had heated it: a lower thermal peak, reached
early, at higher pressure, and a cold path back to the surface.
\end{solution}

\begin{solution}{exo:b2:pt-paths:12}
(a) The crust weighs $35000\times2800\times9.81 = \qty{0.961}{GPa}$; the remaining
\qty{1.839}{GPa} need $1.839\times10^9/(3300\times9.81) = \qty{56.8}{km}$ of mantle: depth
\qty{91.8}{km}, as the \qty{90}{km} or more of the original study. (b) $(4.0 -
0.961)\times10^9/(3300\times9.81) = \qty{93.9}{km}$, depth \qty{128.9}{km}.
(c) $700/91.8 = \qty{7.6}{\celsius\per\km}$: a very high ratio of pressure to
temperature, the series of subduction zones. (d) $(91.8 - 35)/34 = \qty{1.7}{Myr}$.
\end{solution}

\begin{solution}{pb:b2:pt-paths:1}
\noindent\textbf{Part I.}

\textbf{1.} Garnet core $0.080/0.760 = 0.1053$, biotite $0.40/0.58 = 0.6897$; $K_D =
0.1527$, $\ln K_D = -1.8793$.

\textbf{2.} Rim $0.120/0.720 = 0.1667$, matrix biotite $0.46/0.53 = 0.8679$; $K_D =
0.1921$, $\ln K_D = -1.6497$.

\textbf{3.} $3R = 24.942$; numerator $52630 + 0.311\times6000 = 54496$; denominator $20.19
+ 24.942\times1.8793 = 67.064$; $T = \qty{812.6}{K} = \qty{539.5}{\celsius}$.

\textbf{4.} Denominator $20.19 + 24.942\times1.6497 = 61.337$; $T = 54496/61.337 =
\qty{888.5}{K} = \qty{615.3}{\celsius}$.

\textbf{5.} $0.311\times1000/67.064 = \qty{4.6}{K}$ for the core and $0.311\times1000/61.337
= \qty{5.1}{K}$ for the rim.

\textbf{6.} A thermometer applies to minerals that were in equilibrium together. The
biotite included in the core was shut away from the rest of the rock when the core grew;
the matrix biotite last exchanged with the rim. Pairing the core with the matrix would mix
two moments.

\noindent\textbf{Part II.}

\textbf{7.} Core $22.80\times539.5 - 1093 = \qty{11208}{bar}$; rim $22.80\times615.3 -
1093 = \qty{12936}{bar}$.

\textbf{8.} $3\ln(0.050/0.25) = -4.8283$; $RT = 8.314\times812.65 = \qty{6756.4}{J}$; $P =
11208 - 6756.4\times4.8283/6.21 = 11208 - 5253 = \qty{5955}{bar}$.

\textbf{9.} $3\ln(0.060/0.28) = -4.6213$; $RT = 8.314\times888.45 = \qty{7386.6}{J}$; $P
= 12936 - 5497 = \qty{7439}{bar}$.

\textbf{10.} At \qty{7439}{bar}: $T = (52630 + 2314)/61.337 = \qty{895.8}{K} =
\qty{622.6}{\celsius}$; at that temperature $P = 13102 - 5542 = \qty{7560}{bar}$; then $T
= \qty{623.2}{\celsius}$. The rim grew at about \qty{623}{\celsius} and \qty{7.56}{kbar}.

\textbf{11.} At \qty{5954}{bar}, $T = \qty{539.2}{\celsius}$ and $P = \qty{5950}{bar}$: the
core grew at about \qty{539}{\celsius} and \qty{5.95}{kbar}. The first guess,
\qty{6}{kbar}, was already close, and the thermometer moves only \qty{4.6}{K} per kilobar.

\textbf{12.} Core $5950\times10^5/(2800\times9.81) = \qty{21.7}{km}$; rim
$7560\times10^5/27468 = \qty{27.5}{km}$.

\textbf{13.} At \qty{539}{\celsius} the line lies at $3.67 + 0.39\times2.10 =
\qty{4.49}{kbar}$, below \qty{5.95}{kbar}: kyanite, consistent with the kyanite included
in the core. At \qty{623}{\celsius} it lies at $5.77 + 0.23\times2.09 = \qty{6.25}{kbar}$,
below \qty{7.56}{kbar}: kyanite again.

\noindent\textbf{Part III.}

\textbf{14.} $+\qty{84}{\celsius}$ and $+\qty{1.6}{kbar}$ (about \qty{6}{km} deeper): the
rock was being buried and heated, on its \omterm{def:b2:pt-paths:paths}{prograde path}.

\textbf{15.} At \qty{623}{\celsius} or more, below \qty{6.25}{kbar}: after the rim grew the
rock lost at least \qty{1.3}{kbar}, about \qty{5}{km} of overburden, while staying at
least as hot. It was exhumed while still hot.

\textbf{16.} Chlorite is a low-grade, hydrous mineral: it records the late \omterm{def:b2:pt-paths:paths}{retrograde
path}, cooling at shallow depth, with water.

\textbf{17.} Burial and heating to the rim, decompression without cooling, then cooling:
drawn with temperature to the right and pressure upwards, the path turns clockwise. The
pressure peak came first, at the rim; the thermal peak at the same time or later, during
decompression.

\textbf{18.} Yes: the model also reaches its pressure peak first and its thermal peak,
\qty{593}{\celsius} at \qty{7.8}{kbar}, during the ascent, close to the schist's rim
conditions. Thickening followed by slow erosion explains the schist's path.

\noindent\textbf{Part IV.}

\textbf{19.} $(623 - 405)/(412.0 - 398.0) = 218/14.0 = \qty{15.6}{\celsius}$ per million
years.

\textbf{20.} The rock cooled about one and a half times faster than assumed. A mineral
that cools faster has less time to lose its argon, so it closes at a higher temperature:
the true closure temperature lies somewhat above \qty{405}{\celsius}, and the closure
somewhat deeper.

\textbf{21.} $(405 - 10)/30 = \qty{13.2}{km}$.

\textbf{22.} $27.5 - 13.2 = \qty{14.3}{km}$ (with the unrounded depths, 27.52 and 13.17,
\qty{14.35}{km}).

\textbf{23.} $\sqrt{1.5^2 + 1.0^2} = \qty{1.8}{Myr}$: $14.0 \pm \qty{1.8}{Myr}$, so a rate
between $14.35/15.8 = 0.91$ and $14.35/12.2 = \qty{1.18}{km/Myr}$.

\textbf{24.} $14.35/14.0 = \qty{1.03}{km}$ per million years, that is
\textbf{\qty[text-series-to-math=true]{1.0}{mm.yr^{-1}}}, about thirty times slower than
the first stage of \omterm{def:b2:pt-paths:exhumation}{exhumation} of the deepest Alpine rocks, and about twice the erosion rate
of the thickening model.
\end{solution}
